% Three errors from three different corners, and one thing that only looks
% like an error.
\catcode`\{=1 \catcode`\}=2

% §1069 `handle_right_brace`, the `bottom_level` case: a `}` closing a group
% that was never opened. We returned in silence, which is why the help calls it
% harmless — TeX says so and keeps going.
\hbox{a}}

% §977: `\vsplit` of a register holding an \hbox. Two faults here, not one —
% we raised no error, *and* we removed the box before checking its kind, so the
% register was voided. TeX checks first: "I can't split such a box, so I'll
% leave it alone", and the register keeps its contents.
\setbox0\hbox{}
\setbox1\vsplit0 to 5pt
\message{PROBE:kept=\ifvoid0 void\else full\fi}
\message{PROBE:got=\ifvoid1 void\else full\fi}

% ...and a \vbox splits without complaint, so the check is about the kind and
% not about splitting.
\setbox2\vbox{\hrule\vskip10pt\hrule}
\setbox3\vsplit2 to 1pt
\message{PROBE:split=\ifvoid3 void\else full\fi}

% §1296 and §1293 end `end_diagnostic(true); print_err("OK")`, so \showbox and
% \showlists acknowledge through the same routine an error uses — which is why
% `! OK.` carries a context line. §1298 gives it `help0` outside
% \errorstopmode, so no help text follows and it is not counted as an error.
\showboxdepth=2 \showboxbreadth=5
\showbox2
\showlists

% \show and \showthe print their `> …` through §1294/§1297 and stop. No `! OK.`
% follows either, which is the boundary this case pins.
\count1=7
\showthe\count1
% §416 backs up and re-reads a font identifier before returning it.  With
% `\showthe` there is no later scanner to pop that spent input level before
% §311 prints the context, so both the `\font` primitive and a font control
% sequence appear under `<recently read>`.  A family selector reads its number
% afterwards and has already popped the level.
\font\r=cmr10 \r \textfont0=\r
\showthe\font
\showthe\r
\showthe\textfont0
\def\m{x}
\show\m
\end
